% GENERATED FILE. Edit canonical lessons, then run npm run generate:books.
% canonical-source-hash: fnv1a64:57764767fb4c0250
% canonical-lessons: HI-C56-road, HI-C56-village, HI-C56-market, HI-C56-shop, HI-C56-field

\chapter{The Road}
\label{ch:hi-the-road}
\hlchaptermodality{56}

\begin{chapteropening}
\textbf{By the end of this chapter:} \emph{I can name five places along the way when I take my leave.}

Complete HI-C56-field and say all five words in order.
\end{chapteropening}

\section[saṛak]{\hi{सड़क} (saṛak) — a road}

\label{lesson:HI-C56-road}



\begin{quote}
Before the new run of words: what did the last two words of the chapter before mean?
\end{quote}

\subsection*{You'll want to know: \hi{सड़क}}
\textbf{\hi{सड़क}} (\emph{saṛak}) — \textquotedblleft{}a road\textquotedblright{}.

\hi{सड़क} is an everyday word with a disputed history; the derivations offered for it do not agree, and none has become standard.

It sits oddly beside the four words after it in this chapter, which are all traceable. A chapter of five words is a fair sample of a language, and one unsettled origin in five is about the honest rate.

\begin{grammarlens}[title={What you've built}]
The first of five places along the way.
\end{grammarlens}

\subsection*{Your turn}
Say these aloud:

\begin{itemize}
\raggedright
  \item \emph{saṛak}
  \item it once more, slowly
  \item \emph{saṛak}, the thing you set out along
\end{itemize}

\begin{tcolorbox}[breakable,colback=teal!4,colframe=teal!35!black,title={Before you move on}]
What does \hi{सड़क} mean? (\textquotedblleft{}a road\textquotedblright{}.) Say it once more.
\end{tcolorbox}
\section[gā̃v]{\hi{गाँव} (gā̃v) — a village}

\label{lesson:HI-C56-village}



\begin{quote}
What did the word before this one mean?
\end{quote}

\subsection*{You'll want to know: \hi{गाँव}}
\textbf{\hi{गाँव}} (\emph{gā̃v}) — \textquotedblleft{}a village\textquotedblright{}.

From Sanskrit \hi{ग्राम} (\emph{grāma}), 'village'. The \emph{gr-} cluster simplified and the \emph{m} left its nasality behind in the vowel as it went — which is exactly what you watched happen to \hi{पाँच} when you first counted to five.

\emph{Grāma} has not disappeared. It is still there, unworn, in the names of thousands of places and in the formal register, so the old form and the new one sit on the same map.

\begin{grammarlens}[title={What you've built}]
The second of five, and the same nasal trade as \hi{पाँच}.
\end{grammarlens}

\subsection*{Your turn}
Say these aloud:

\begin{itemize}
\raggedright
  \item \emph{gā̃v}
  \item it once more, slowly
  \item \emph{gā̃v}, and feel the nasal where the \emph{m} used to be
\end{itemize}

\begin{tcolorbox}[breakable,colback=teal!4,colframe=teal!35!black,title={Before you move on}]
What does \hi{गाँव} mean? (\textquotedblleft{}a village\textquotedblright{}.) Say it once more.
\end{tcolorbox}
\section[bāzār]{\hi{बाज़ार} (bāzār) — a market}

\label{lesson:HI-C56-market}



\begin{quote}
Name the 2 words of this chapter so far.
\end{quote}

\subsection*{You'll want to know: \hi{बाज़ार}}
\textbf{\hi{बाज़ार}} (\emph{bāzār}) — \textquotedblleft{}a market\textquotedblright{}.

Persian \emph{bāzār}, a market. English borrowed the same word — \emph{bazaar} — from the same Persian source, by way of Italian traders.

So this is not a cousin but a shared borrowing: Hindi and English both took it from Persian, separately. You met that shape with \hi{चाय}, which the world took from Chinese by two different roads.

\begin{grammarlens}[title={What you've built}]
The third of five, and a word English took from the same place.
\end{grammarlens}

\subsection*{Your turn}
Say these aloud:

\begin{itemize}
\raggedright
  \item \emph{bāzār}
  \item it once more, slowly
  \item \emph{bāzār}, on the \emph{saṛak}
\end{itemize}

\begin{tcolorbox}[breakable,colback=teal!4,colframe=teal!35!black,title={Before you move on}]
What does \hi{बाज़ार} mean? (\textquotedblleft{}a market\textquotedblright{}.) Say it once more.
\end{tcolorbox}
\section[dukān]{\hi{दुकान} (dukān) — a shop}

\label{lesson:HI-C56-shop}



\begin{quote}
Name the 3 words of this chapter so far.
\end{quote}

\subsection*{You'll want to know: \hi{दुकान}}
\textbf{\hi{दुकान}} (\emph{dukān}) — \textquotedblleft{}a shop\textquotedblright{}.

An Arabic word, \emph{dukkān}, meaning a shop or the raised bench a shopkeeper sits on, which reached Hindi through Persian.

That is the same road that brought you \hi{शुक्रिया} and \hi{मुलाक़ात}: Arabic word, Persian carrier, Hindi destination. By now the route should be as familiar as the words on it.

\begin{grammarlens}[title={What you've built}]
The fourth of five, and a well-travelled road you have walked before.
\end{grammarlens}

\subsection*{Your turn}
Say these aloud:

\begin{itemize}
\raggedright
  \item \emph{dukān}
  \item it once more, slowly
  \item \emph{dukān}, in the \emph{bāzār}
\end{itemize}

\begin{tcolorbox}[breakable,colback=teal!4,colframe=teal!35!black,title={Before you move on}]
What does \hi{दुकान} mean? (\textquotedblleft{}a shop\textquotedblright{}.) Say it once more.
\end{tcolorbox}
\section[khet]{\hi{खेत} (khet) — a field}

\label{lesson:HI-C56-field}



\begin{quote}
Name the 4 words of this chapter so far.
\end{quote}

\subsection*{You'll want to know: \hi{खेत}}
\textbf{\hi{खेत}} (\emph{khet}) — \textquotedblleft{}a field\textquotedblright{}.

From Sanskrit \hi{क्षेत्र} (\emph{kṣetra}), 'field, ground'. The \emph{kṣ-} cluster wore down to \emph{kh-}, and what is left is one short syllable.

\hi{क्षेत्र} survives beside it, intact, meaning 'a region, a domain, a field of study' — the abstract sense. One ancestor, two descendants: the dirt and the discipline.

\begin{grammarlens}[title={What you've built}]
The fifth of five. That is the way out of the village.
\end{grammarlens}

\subsection*{Your turn}
Say these aloud:

\begin{itemize}
\raggedright
  \item \emph{khet}
  \item it once more, slowly
  \item \emph{khet}, \emph{gā̃v}, \emph{saṛak}, \emph{bāzār}, \emph{dukān} — walk it in order
\end{itemize}

\begin{tcolorbox}[breakable,colback=teal!4,colframe=teal!35!black,title={Before you move on}]
What does \hi{खेत} mean? (\textquotedblleft{}a field\textquotedblright{}.) Say it once more.
\end{tcolorbox}
